\documentclass[10pt, aspectratio=169]{beamer}

% --- TEMA Y ESTILO ---
\usetheme{Madrid}
\usecolortheme{whale}

% --- PAQUETES ---
\usepackage[utf8]{inputenc}
\usepackage[spanish]{babel}
\usepackage{amsmath, amssymb}
\usepackage{siunitx}
\usepackage{booktabs}
\usepackage{tabularx}
\usepackage[american]{circuitikz}
\usepackage{pgfplots}
\pgfplotsset{compat=1.18}

\sisetup{output-decimal-marker = {,}, per-mode = symbol}

% --- INFORMACIÓN DEL CURSO ---
\title[Rectificación y Filtros]{Rectificación y Filtrado Capacitivo}
\subtitle{Conversión CA-CD, Dimensionamiento de Rizado y Esfuerzos}
\author{M.Sc. Ing. Bernardo Quiroga Turdera}
\institute[UCB Tarija]{Circuitos Electrónicos III (IMT-221)}
\date{Semana 3: Sesión 5}

\begin{document}

\begin{frame}
  \titlepage
\end{frame}

\begin{frame}{Auditoría Socrática de Entrada}
  \begin{block}{Conectando con la Sesión 4}
    En la clase anterior, analizamos que para señales muy pequeñas (audio, ruido), el diodo se comporta como una resistencia matemática:
    \begin{equation*}
        r_d \approx \frac{\eta V_T}{I_{DQ}}
    \end{equation*}
  \end{block}
  \vspace{0.3cm}
  \textbf{La Paradoja de hoy:}\\
  ¿Por qué en una etapa de rectificación de potencia para transformar $\qty{220}{\volt}$ de AC a DC, ese modelo lineal ($r_d$) es \textbf{completamente inútil} y debemos usar el modelo de caída constante ($V_\gamma = \qty{0.7}{\volt}$) o Lineal por Tramos (PWL)?
\end{frame}

\begin{frame}{El Reto en el PFI (Contexto de Ingeniería)}
  \begin{itemize}
      \item Todo actuador mecatrónico o biosensores requiere un bus de Corriente Continua (DC) muy estable a partir de la red.
      \item \textbf{El Desafío:} La rectificación sin filtro entrega una tensión pulsante que hace colapsar a los microcontroladores y amplificadores (Hito 2).
  \end{itemize}
  
  \vspace{0.3cm}
  \begin{block}{Objetivo de Diseño}
      Diseñar una red de rectificación y filtrado que \textbf{minimice el rizado residual ($V_r$)} garantizando que NO destruiremos los diodos por el infame \textit{surge current} (sobrecorriente repetitiva).
  \end{block}
\end{frame}

\begin{frame}{Comparativa de Topologías de Rectificación}
  \begin{table}
    \centering
    \renewcommand{\arraystretch}{1.5}
    \begin{tabularx}{\textwidth}{@{} >{\bfseries}l X X X @{}}
      \toprule
      Parámetro de Diseño & Media Onda & Onda Completa (Tap) & \textcolor{blue}{Puente de Graetz} \\
      \midrule
      Diodos requeridos & 1 & 2 & \textcolor{blue}{4} \\
      Tensión pico salida & $V_p - V_\gamma$ & $V_p - V_\gamma$ & \textcolor{blue}{$V_p - 2V_\gamma$} \\
      Frecuencia rizado ($f_r$) & $f_{in}$ (\qty{50}{\hertz}) & $2f_{in}$ (\qty{100}{\hertz}) & \textcolor{blue}{$2f_{in}$ (\qty{100}{\hertz})} \\
      Tensión Inversa (PIV) & $2V_p$ & $2V_p$ & \textcolor{blue}{$V_p$} \\
      Eficiencia máxima & $40.6\%$ & $81.2\%$ & \textcolor{blue}{$81.2\%$} \\
      \bottomrule
    \end{tabularx}
  \end{table}
\end{frame}

\begin{frame}{Mecánica del Filtrado Capacitivo}
  \begin{columns}
      \begin{column}{0.4\textwidth}
          \begin{itemize}
              \item \textbf{Carga:} El diodo conduce y eleva $C$ hasta $V_{p(\text{out})}$.
              \item \textbf{Corte:} La fuente cae. El diodo se polariza en inversa.
              \item \textbf{Descarga:} El capacitor asume toda la carga del circuito mediante:
              $$v_o(t) = V_{p} e^{-\frac{t}{R_L C}}$$
          \end{itemize}
      \end{column}
      \begin{column}{0.6\textwidth}
          \begin{tikzpicture}[scale=0.7]
              \begin{axis}[
                  axis lines=middle, xmin=0, xmax=10, ymin=0, ymax=1.2,
                  xtick={1.57, 4.71}, xticklabels={$\frac{T}{4}$, $\frac{3T}{4}$},
                  ytick={1, 0.7}, yticklabels={$V_{p}$, $V_p - V_r$},
                  grid=both, grid style={dashed, gray!30}
              ]
              \addplot[domain=0:10, samples=150, gray, dashed, thick] {abs(sin(deg(x)))};
              \addplot[domain=0:1.57, samples=30, blue, ultra thick] {sin(deg(x))};
              \addplot[domain=1.57:3.7, samples=30, blue, ultra thick] {exp(-0.15*(x-1.57))};
              \addplot[domain=3.7:4.71, samples=30, blue, ultra thick] {abs(sin(deg(x)))};
              \addplot[domain=4.71:6.84, samples=30, blue, ultra thick] {exp(-0.15*(x-4.71))};
              \addplot[domain=6.84:7.85, samples=30, blue, ultra thick] {abs(sin(deg(x)))};
              \draw[<->, red, thick] (axis cs: 6.84, 0.72) -- (axis cs: 6.84, 1) node[midway, right] {$V_r$};
              \end{axis}
          \end{tikzpicture}
      \end{column}
  \end{columns}
\end{frame}

\begin{frame}{Formulación Analítica: Rizado y Tensión Continua}
  \begin{block}{1. Tensión de Rizado Pico a Pico ($V_r$)}
      Aproximando la descarga exponencial a una recta (si $\tau \gg T$):
      \begin{equation*}
          V_r \approx \frac{I_L}{2 f_{in} C} \quad (\text{Onda Completa})
      \end{equation*}
  \end{block}
  
  \vspace{0.3cm}
  \begin{block}{2. Tensión Media en Continua ($V_{DC}$)}
      Es el valor útil real que "ven" los reguladores posteriores:
      \begin{equation*}
          V_{DC} = V_{p(\text{out})} - \frac{V_r}{2}
      \end{equation*}
  \end{block}
\end{frame}

\begin{frame}{El Costo Oculto: Corriente Pico en el Diodo}
  \alert{\textbf{Ley de Conservación de Energía:}}\\
  A mayor capacitancia $C$, \textbf{menor} es el rizado $V_r$, pero \textbf{menor} es el tiempo de conducción ($\Delta t$). Toda la carga extraída se debe reponer en un instante.
  
  \vspace{0.2cm}
  \textbf{Ángulo de conducción:} $\theta_c \approx \sqrt{\frac{2 V_r}{V_p}}$
  
  \vspace{0.2cm}
  \begin{block}{Corriente Pico Repetitiva ($I_{D(\max)}$)}
      \begin{equation*}
          I_{D(\max)} \approx I_L \left( 1 + \pi \sqrt{\frac{2 V_{p(\text{out})}}{V_r}} \right)
      \end{equation*}
  \end{block}
  \textit{Riesgo Técnico:} Un capacitor sobredimensionado puede destruir la unión P-N del diodo en los primeros ciclos por calor disipado en $r_D$.
\end{frame}

\begin{frame}{Resolución en Pizarra: Fuente PFI}
  \begin{columns}
      \begin{column}{0.4\textwidth}
          \begin{circuitikz}[american, scale=0.6, transform shape]
              \draw (-2,0) to[sV, l=\qty{12}{\volt_{\text{rms}}}] (-2,-4);
              \draw (-2,0) |- (0,0); \draw (-2,-4) |- (4,-4) -- (4,0);
              \draw (0,0) to[D, *-*] (2,2); \draw (4,0) to[D, *-*] (2,2);
              \draw (2,-2) to[D, *-*] (0,0); \draw (2,-2) to[D, *-*] (4,0);
              \draw (2,2) -- (5,2) to[C] (5,-2) -- (2,-2);
              \draw (5,2) -- (7,2) to[R, l=\qty{500}{\milli\ampere}] (7,-2) -- (5,-2);
          \end{circuitikz}
      \end{column}
      
      \begin{column}{0.6\textwidth}
          \textbf{Resultados Analíticos ($V_r \le \qty{0.5}{\volt}$):}
          \begin{itemize}
              \item $V_{p(\text{out})} = 12\sqrt{2} - 1.4 = \mathbf{\qty{15.57}{\volt}}$
              \item $C_{\text{req}} = \frac{0.5}{100 \cdot 0.5} = \mathbf{10\,000\,\mu\text{F}}$
              \item $I_{D(\max)} = 0.5 \left( 1 + \pi \sqrt{\frac{31.14}{0.5}} \right) = \mathbf{\qty{12.89}{\ampere}}$
              \item PIV = $16.97 - 0.7 = \mathbf{\qty{16.27}{\volt}}$
          \end{itemize}
          \alert{\textbf{Decisión:}} Diodo 1N4148 se quema. Usaremos \textbf{1N4007}.
      \end{column}
  \end{columns}
\end{frame}

\begin{frame}{Taller Computacional (Python)}
  \begin{center}
      ¡Es hora de validar nuestro diseño matemático!
  \end{center}
  \vspace{0.2cm}
  En sus laboratorios, ejecutarán el \textit{script} numérico que simula la carga/descarga exponencial ciclo a ciclo.
  \vspace{0.4cm}
  
  \textbf{Reto de Auditoría IA:}\\
  ¿Qué pasa físicamente con el pico de corriente $I_{D(\max)}$ si bajamos deliberadamente el capacitor a $100\,\mu\text{F}$ permitiendo mucho ruido en el bus DC? \\
  \textit{Analicen la forma de onda del Subplot 2.}
\end{frame}

\end{document}